En somme, que faut-il retenir du rôle des bibliothèques musicales, et de leurs plateformes numériques, dans la patrimonialisation et la valorisation du patrimoine musical de la pratique historiquement informée ?

Comme nous l'avons vu, les bibliothèques musicales font figure d'exception dans le paysage des bibliothèques. Les spécificités de leurs collections, ainsi que la complexité de la classification de leurs ressources selon les typologies traditionnelles des archives, en font des terrains délicats pour l'application des normes bibliothéconomiques et archivistiques. Associée à cette difficulté technique, s'ajoute une difficulté métier puisque le personnel n'est pas toujours spécialisé en documentation : souvent administrateurs ou musiciens, les bibliothécaires de ces structures n'ont parfois de bibliothécaire que le nom. 


Les bibliothèques musicales sont donc des institutions difficiles à catégoriser, tant les réalités du terrain sont singulières. Ce ne sont pas des bibliothèques de conservation, mais elles n'ont pas non plus vocation à être des bibliothèques de lecture publique. A l'exception de la Médiathèque Musicale de Paris dont le modèle est très singulier, s'agissant d'une médiathèque spécialisée dans la musique, les autres bibliothèques musicales n'ont pas pour vocation première l'accueil du public. Ce sont majoritairement des bibliothèques de recherche destinées aux musicologues et musiciens. 
Les bibliothèques d'orchestre forment quant à elles une catégorie à part au sein même des bibliothèques musicales. Organes de fonctionnement interne des orchestres, elles ont un fonctionnement unique, entre conservation et \textit{records management}. 


A première vue, les bibliothèques d'orchestre paraissent être des institutions de conservation car leurs fonds ont vocation à documenter les phases de production des représentations de l'orchestre. Dans la plateforme LaDocumenta.eu, ce circuit de conservation du document en fonction de la phase de production à laquelle il appartient - pré-production, production, concert ou post production - illustre bien cette logique.


Mais en analysant attentivement ces institutions, il peut être plus juste d'y voir un organisme de \textit{records management} car leurs documents peuvent être assimilés à ce qui s'appelle, dans les pays anglo-saxons, les \textit{records}, pouvant se traduire par \enquote{documents d'activité} même si la traduction de ce terme ne fait pas consensus\footnote{Cette réflexion sur la catégorisation des documents des bibliothèques d'orchestre a utilisé à profit le cours d'\enquote{Archivistique contemporaine} dispensé en première année de master \enquote{Technologies numériques appliquées à l'histoire} (2024-2025) par Edouard Vasseur, directeur d'études à l'Ecole des chartes - titulaire de la chaire Histoire des institutions, diplomatique et archivistique contemporaine - et sa bibliographie. La notion de \textit{records} et sa traduction est notamment issue de l'article de Sylvie Dessolin Baumann dans le numéro 228 de La Gazette des archives dont la référence bibliographique est : Dessolin Baumann (Sylvie),\textit{ Des chartriers aux bases de données. Les enjeux de la gestion des documents d’activité.} Dans La Gazette des archives, n°228, 2012-4. Normalisation et gestion des documents d’activité (records management) : enjeux et nouvelles pratiques pour notre profession. pp. 93-102. DOI : \url{https://doi.org/10.3406/gazar.2012.4987}}. 


En effet, les documents musicaux des bibliothèques d'orchestre, s'ils sont patrimonialisés pour les annotations qu'ils contiennent, très précieuses pour la pratique historiquement informée, restent pourtant des archives courantes, utilisés par l'orchestre pour ses productions musicales\footnote{Sur la spécificité des collections des bibliothèques musicales : Hausfater (Dominique) et Campos (Rémy), \textit{L’interprétation musicale dans les fonds des bibliothèques }: Journée d’étude du 7 Janvier 2006, Conservatoire de Musique de Genève / Haute École de Musique. Fontes Artis Musicae, vol. 54, no.1, 2007, pages.1–5. JSTOR, \url{http://www.jstor.org/stable/23510565}.(consulté le 25 août 2026)}. 

Ainsi cette double vocation est au coeur de la mission des bibliothèques d'orchestre qui doivent composer avec des exigences différentes : favoriser les recherches musicologiques sur les fonds conservés mais également conserver durablement ces fonds patrimoniaux qui, orientés vers les représentations, documentent les choix interpétatifs. Ce rôle prend tout son sens dans le cadre de pratiques historiquement informées de la musique, où les chefs sont particulièrement attentifs à la tradition d'interprétation des oeuvres. 


La bibliothèque d'orchestre est donc un lieu clé tant pour les musiciens et chefs qui construisent leur interprétation que pour les chercheurs en musicologie. Avec les enjeux de science ouverte, fondamentaux dans les bibliothèques universitaires aujourd'hui, la fonction clé des bibliothèques musicales est la mise en relation des acteurs de la recherche musicologique avec les musiciens de pratique historiquement informée particulièrement sensibles aux \textit{performance studies}. C'est précisément la fonction recherchée par le consortium d'orchestres européens LaDocumenta.eu, dont l'objectif est de mettre en relation les acteurs de la pratique historiquement informée, qu'ils soient chercheurs, musiciens ou encore facteurs d'instruments et pédagogues.

Le grand public, s'il n'est pas l'usager attendu des bibliothèques ne doit pas pour autant être exclu de ces initiatives. C'est le pari pris par LaDocumenta.eu qui compte le grand public dans ses objectifs d'audience de la plateforme.  Même si le grand public n'a pas vocation à fréquenter les bibliothèques des orchestres du consortium en dehors d'événements dédiés, celui-ci fait néanmoins l'objet de contenus ciblés dans la plateforme. 

Que serait la pratique historiquement informée sans son public ? Afin de faire connaître ces plateformes et de conformer ces bibliothèques singulières à la vocation première d'ouverture des bibliothèques, ce type de plateforme gagne à inclure le grand public ne serait-ce que dans ses actions de médiation, comme dans des initiatives de navigation spécialisée, avec plusieurs niveaux de lecture. 


Les outils d'intelligence artificielle, assistants documentaires capables d'épauler les agents dans leurs activités répétitives de dépôt de ressources dans des bibliothèques numériques, permettent aussi la mise en oeuvre de ces navigations assistées pour les usagers des plateformes. A la fois assistants de recherche pour les chercheurs, coachs personnalisés pour les musiciens et guides de la plateforme pour le grand public, les outils d'IA sont prometteurs pour les nouvelles mutations des métiers et pratiques des bibliothèques. 


Toutefois, si ces outils prometteurs semblent être une évidence dans ce domaine où les grandes masses documentaires sont un \enquote{terreau favorable\footnote{Moreux (Jean-Philippe), \textit{Indexation automatique des documents : de nouvelles métadonnées pour de nouveaux services.} Dans Cavalié (Étienne) (dir.), \textit{L'indexation matière en transition : de la réforme Rameau à l'indexation automatique.} Paris, Éditions du Cercle de la librairie, coll. « Bibliothèques », 2019, p. 113-133.}} à l'utilisation de l'IA, ces promesses pour l'avenir ne sont pas exemptes de questionnements, d'enjeux et de limites qui se manifestent par des contraintes bibliothéconomiques, techniques, juridiques et éthiques pour les bibliothèques. 

En effet, pour pouvoir intégrer l'IA dans leurs pratiques, les bibliothèques doivent transformer les ressources documentaires, parfois hétérogènes dans leur qualité et leur description, en données structurées et interopérables afin qu'elles puissent être exploitables par l'IA et réutilisables, hors des silos applicatifs. 


Les plateformes numériques musicales comme LaDocumenta.eu offrent, comme nous l'avons vu, des solutions originales pour transformer le patrimoine des bibliothèques d'orchestre en données de recherche exploitables. 

La transformation technique, en données structurées grâce aux référentiels du web sémantique, et la réflexion intellectuelle quant à l'utilisation du document et son organisation, permettent de convertir ces ressources dans une logique d'interopérabilité des données. Ces évolutions s'inscrivent pleinement dans la dynamique de la Transition bibliographique et la mise en œuvre du modèle IFLA LRM. 
Ainsi, les ressources documentaires musicales transformées sont des \enquote{terrains d'expérimentation\footnote{Carbone (Pierre), \textit{Chapitre XII : Les services en ligne}. Dans : \textit{Les bibliothèques}. Paris: Presses Universitaires de France. Que sais-je ?, 2017, p. 120.}} favorables à l'utilisation de l'IA. 

Dans cette ère numérique, la politique documentaire traditionnelle se double d'une politique de gouvernance des données, indispensable pour envisager, un jour, l'intégration de l'IA dans les tâches documentaires. Au sein de ce terreau fertile, les outils d'IA peuvent alors devenir des assistants documentaires pertinents nécessitant toutefois une vigilance dans l'organisation des tâches. L'IA ne doit en effet pas être utilisée à des fins de facilité : ces éléments, s'ils doivent être pris en compte et favorisés dans le travail quotidien, ne peuvent devenir des mots d'ordre qui s'exercent aux dépens de la rigueur bibliothéconomique.. Au contraire, elle doit être vue comme un outil permettant aux équipes de garder un temps précieux pour des tâches de valorisation des collections, de médiation, de conservation préventive, qui sont parfois remplacées par des tâches répétitives de catalogage notamment. En aucun cas l'IA ne doit être utilisée pour elle-même, comme pour se conformer à une mode de l'intelligence artificielle. A l'heure où les campagnes culturelles de financement de la recherche semblent encourager l'utilisation de l'IA, ces utilisations doivent être plus que jamais délimitées, précisées et critiquées. 

Contrairement aux idées reçues, ces campagnes de financement n'encouragent pas réellement l'utilisation de l'IA mais incitent les institutions à se saisir de ces opportunités pour repenser leurs pratiques et leurs fonctionnements, dans une démarche raisonnée. A l'image de la BnF qui a élaboré une feuille de route de l'utilisation de l'IA, les institutions doivent se documenter sur ces outils afin d'adopter une stratégie raisonnée de leur utilisation. 

En effet, bien que prometteurs, ces outils donnant l'illusion d'une simplicité d'utilisation et de simplification du travail quotidien présentent des limites et des contraintes d'utilisation, en raison de considérations juridiques, éthiques et techniques qui imposent un management stratégique conscient de ces enjeux. 


L'aspect éthique de l'IA a fait couler beaucoup d'encre, tant dans les publications scientifiques que dans les revendications syndicales. En bibliothèques, l'IA n'a pas vocation à supprimer de l'emploi, comme cela peut être parfois entendu, mais vise à aider les agents face à la multiplication des masses documentaires. Vus comme assistants documentaires ou guides de navigation des plateformes numériques, ces outils sont censés être utilisés comme support et non comme remplacement des personnels des bibliothèques. Plus encore, ce doivent être les bibliothèques qui mènent l'innovation dans ce domaine : dès lors que les bornes de l'intelligence artificielle sont fermement définies et encadrées par les institutions, son usage est adapté aux besoins et ne comporte, en théorie, aucun risque pour la sûreté des collections ou en terme de la concurrence à l'égard des personnels.




En replaçant l’humain au cœur des opérations de médiation et d’analyse scientifique, ces transformations numériques ne dénaturent pas les bibliothèques d’orchestre : elles révèlent la richesse d’un patrimoine musical vivant, aujourd’hui accessible, exploitable et réutilisable. L’intégration de ces outils d’intelligence artificielle au sein de plateformes musicales numériques à l’instar de LaDocumenta.eu dépasse alors les enjeux techniques précités. Plus de trente ans après le succès du film Tous les matins du monde, la médiation du patrimoine musical et l’action culturelle musicale en bibliothèque ne passent plus seulement par l’extase des concerts ou du cinéma, mais également par des plateformes numériques capables de faire le lien entre la source patrimoniale, la scène et son public. 

L’intégration de ces outils dans les pratiques actuelles des bibliothèques suppose, il est vrai, d’importants questionnements face aux transformations du métier et à l’accompagnement au changement qui doit en découler. Toutefois, ces outils redéfinissent le rôle du bibliothécaire musical qui, toujours tourné vers la conservation patrimoniale, devient aussi un professionnel de la donnée attentif à la mémoire de l’interprétation musicale.

Réfléchir à la patrimonialisation d’un mouvement musical ancre la pratique historiquement informée dans le spectacle vivant immatériel. Si la patrimonialisation des partitions est un champ parfaitement maîtrisé des institutions patrimoniales aujourd’hui, la patrimonialisation des choix interprétatifs des ensembles de pratique historiquement informée, à l’aune de l’intelligence artificielle et des enjeux d'interopérabilité dans un contexte Transition bibliographique, s’inscrit pleinement dans les débats internationaux actuels. 
Les travaux de la communauté AI4LAM (Artificial Intelligence for Libraries, Archives and Museums)  et de sa conférence annuelle Fantastic Futures l’illustrent. 
Ces conférences, en réunissant institutions patrimoniales, chercheurs et spécialistes de la donnée témoignent de cette volonté des institutions d’organiser la définition des cadres éthiques, juridiques et techniques de ces outils, au service de la recherche et de la création artistique.
